\documentclass[conference]{IEEEtran}
\usepackage[utf8]{inputenc}
\usepackage[T1]{fontenc}
\usepackage[spanish]{babel}

\begin{document}

\title{Covergroups en Verificación Funcional\\ \large Tecnológico de Costa Rica - Curso de Verificación Funcional}
\author{Fabián Chacón Solano - 2018135154}
\maketitle

\section{¿Qué es un covergroup en verificación funcional?}

Un \textit{covergroup} es una construcción de SystemVerilog utilizada para
medir cobertura funcional. Su propósito es registrar qué valores, eventos o
combinaciones relevantes del diseño han sido ejecutados durante una
simulación. A diferencia de la cobertura de código, que indica qué partes de la
implementación fueron ejecutadas, la cobertura funcional está relacionada con
la intención del diseño y con el plan de verificación \cite{Spear2008}, por ende
es una herramienta clave para evaluar el progreso de la verificación y asegurar que se e
están evaluando los escenarios.

En otras palabras, un \textit{covergroup} permite responder preguntas como:
``¿se probaron todos los tipos de operación?'', ``¿se ejercitaron los casos
borde?'' o ``¿se observaron las combinaciones importantes de entradas?''.
Spear explica que la cobertura funcional ayuda a medir el progreso de la
verificación y a encontrar huecos de cobertura que requieren nuevos estímulos
o ajustes en las pruebas \cite{Spear2008}, este trabajo es fundamental
puesto que humanamente es considerablemente complejo cubrir todos los
escenarios posibles, por lo que el uso de covergroups es una herramienta relevante
para el diseño funcional de circuitos integrados.

\section{¿Qué partes tiene un covergroup?}

Un \textit{covergroup} está formado principalmente por los siguientes elementos:

\begin{itemize}
    \item \textbf{Covergroup}: bloque que agrupa una o más mediciones de
    cobertura. Puede ser muestreado explícitamente con el método
    \texttt{sample()} o mediante un evento declarado en su definición.

    \item \textbf{Coverpoints}: variables o expresiones que se desean observar.
    Cada \textit{coverpoint} mide qué valores de una señal, variable o expresión
    fueron ejecutados durante la simulación \cite{Spear2008}.

    \item \textbf{Bins}: contadores asociados a valores o rangos de valores.
    Los \textit{bins} son la unidad básica de medición de cobertura funcional,
    cuando se muestrea un valor que cae dentro de un \textit{bin}, ese bin queda
    cubierto \cite{Spear2008}.

    \item \textbf{Cross coverage}: medición de combinaciones entre dos o más
    \textit{coverpoints}. Se usa cuando no basta con saber que dos valores
    aparecieron por separado, sino que es requereido para comprobar que aparecieron en
    combinación \cite{Spear2008}.

    \item \textbf{Opciones de cobertura}: parámetros como
    \texttt{option.per\_instance}, \texttt{option.goal},
    \texttt{option.weight} o límites para bins automáticos. Estas opciones
    permiten controlar cómo se calcula y reporta la cobertura \cite{Spear2008}.

    \item \textbf{Bins especiales}: SystemVerilog permite usar
    \texttt{ignore\_bins} para excluir valores que no deben afectar el cálculo,
    también \texttt{illegal\_bins} para marcar valores que representan situaciones
    invalidas \cite{Spear2008}.
\end{itemize}

\section{¿Por qué es importante la implementación de covergroups?}

La implementación de \textit{covergroups} es importante porque permite medir
de manera explícita si el testbench realmente está verificando los escenarios
definidos en el plan de verificación. Sin cobertura funcional, una simulación
puede terminar sin errores y aun así no haber probado casos importantes del
diseño.

Spear resalta que la cobertura funcional se usa dentro de un ciclo de
retroalimentación, osea se ejecutan pruebas, se analiza la cobertura alcanzada, se
identifican huecos y luego se modifican los estímulos o restricciones para
cubrir escenarios faltantes \cite{Spear2008}. Esto permite mejorar la calidad
del proceso de verificación y tomar decisiones basadas en datos, no solo en la
cantidad de pruebas ejecutadas.

Además, los \textit{covergroups} ayudan a diferenciar entre haber ejecutado
código y haber verificado funcionalidad. Un diseño puede alcanzar alta
cobertura de código, pero seguir sin probar ciertas condiciones funcionales
importantes. Por eso, la cobertura funcional complementa la cobertura de
código y permite evaluar el avance frente a la especificación \cite{Spear2008}.

\section{Recomendaciones para la creación de covergroups}

Para crear \textit{covergroups} útiles se recomiendan las siguientes prácticas:

\begin{itemize}
    \item \textbf{Partir del plan de verificación}. Los \textit{coverpoints}
    deben representar comportamientos, valores y combinaciones importantes
    para la especificación, no señales escogidas al azar.

    \item \textbf{Definir bins significativos}. Para señales grandes, no es conveniente
    depender siempre de bins automáticos. Es preferible definir rangos útiles,
    casos borde y valores especiales, como ceros, máximos, mínimos o estados
    relevantes \cite{Spear2008}.

    \item \textbf{Evitar cobertura decorativa}. Un \textit{covergroup} debe medir
    algo que ayude a tomar decisiones de verificación. Si una métrica no cambia
    la estrategia de pruebas, probablemente no aporta valor.

    \item \textbf{Muestrear en el momento correcto}. La cobertura debe
    muestrearse cuando los datos sean válidos, por ejemplo al completarse una
    transacción o cuando el monitor observa una salida estable. Spear indica
    que las dos partes principales de la cobertura funcional son los datos
    muestreados y el momento en que se muestrean \cite{Spear2008}.

    \item \textbf{Deshabilitar cobertura durante reset}. Se recomienda evitar
    que valores de inicialización o reset contaminen la cobertura. Esto puede
    hacerse con condiciones como \texttt{iff (!reset)} \cite{Spear2008}.

    \item \textbf{Usar cross coverage con cuidado}. La cobertura cruzada puede
    crecer rápidamente, ya que el número de bins aumenta con las combinaciones.
    Por eso debe usarse solo cuando la relación entre variables sea realmente
    importante \cite{Spear2008}.

    \item \textbf{Usar \texttt{ignore\_bins} e \texttt{illegal\_bins} cuando
    corresponda}. Los valores imposibles o irrelevantes deben excluirse del
    cálculo, mientras que los valores inválidos deben marcarse como ilegales.

    \item \textbf{Analizar los reportes de cobertura}. La cobertura no debe ser
    solo un número final. Deben revisarse los bins sin cubrir para decidir si
    se necesitan más datos aleatorios, pruebas dirigidas o cambios en las
    restricciones.
\end{itemize}

\begin{thebibliography}{1}

\bibitem{Spear2008}
C. Spear, \textit{SystemVerilog for Verification: A Guide to Learning the
Testbench Language Features}. New York, NY, USA: Springer Science+Business
Media, 2008.

\end{thebibliography}

\end{document}



